\documentclass[11pt]{article}
\usepackage[paperwidth=220mm,paperheight=320mm,left=16mm,right=16mm,top=18mm,bottom=18mm,headheight=16pt,headsep=10pt,footskip=14pt]{geometry}
\usepackage{fix-cm}
\usepackage{fontspec,amsmath,amssymb,graphicx,array,multirow,tikz}
\usetikzlibrary{decorations.pathreplacing}
\usepackage[unicode,hidelinks]{hyperref}
\usepackage{polyglossia}
\setmainlanguage{latin}
\setotherlanguages{arabic,greek,hebrew,syriac}
\setmainfont{Linux Libertine O}[Ligatures=TeX]
\setsansfont{Linux Biolinum O}[Ligatures=TeX]
\newfontfamily\greekfont{FreeSerif}[Script=Greek]
\newfontfamily\arabicfont{Amiri}[Script=Arabic]
\newfontfamily\hebrewfont{Frank Ruhl Hofshi}[Script=Hebrew]
\newfontfamily\syriacfont{Segoe UI Historic}[Script=Syriac]
\newfontfamily\CapFont{Noto Serif}
\newfontfamily\MoonFont{FreeSerif}
\newcommand{\MoonSym}{\text{\MoonFont ☾}}\newcommand{\SunSym}{\text{\MoonFont ⊙}}
\newfontfamily\EthiopicFont{Ebrima}
\newcommand{\textethiopic}[1]{{\EthiopicFont #1}}
\newfontfamily\NallinoSigns{NallinoSigns.otf}[Path=./fonts/]
\newcommand{\AbjadZero}{{\NallinoSigns\symbol{"E000}}}
\hypersetup{pdftitle={Al-Battani, Opus astronomicum. Nallino, versio Latina et adnotationes (Pars I). S02},pdfauthor={Al-Battani; Carlo Alfonso Nallino (Latin translation and notes)}}
\makeatletter
\def\ps@sourceedition{%
\def\@oddhead{\vbox{\hbox to\textwidth{\fontsize{8}{10}\selectfont C. A. NALLINO — VERSIO ET ADNOTATIONES\hfil\leftmark}\vskip3pt\hrule height.25pt}}%
\let\@evenhead\@oddhead
\def\@oddfoot{\hbox to\textwidth{\fontsize{7}{9}\selectfont\rightmark\hfil\thepage}}%
\let\@evenfoot\@oddfoot}
\makeatother
\pagestyle{sourceedition}
\setlength{\parindent}{0pt}\setlength{\parskip}{0pt}
\newcommand{\BeginSource}[2]{\clearpage\markboth{PAG. #2}{#1}\hypertarget{#1}{}\pdfbookmark[0]{#1 — p. #2}{bk-#1}\fontsize{11.1}{13.8}\selectfont}
\tracinglostchars=2
\newcommand{\Name}[1]{{\addfontfeatures{LetterSpace=4.0}#1}}
\newcommand{\Indent}{\hspace*{1.6em}}
% an Arabic run inside a Latin line: its descenders do not open the line (the printed lines are evenly spaced)
\newcommand{\ArabicRun}[1]{\smash[b]{\textarabic{#1}}}
\newcommand{\qfrac}[2]{{}^{#1}\!/_{{#2}}}
\newcommand{\sa}{\textsuperscript{\textit{a}}}\newcommand{\sm}{\textsuperscript{\textit{m}}}\newcommand{\sd}{\textsuperscript{\textit{d}}}
\newcommand{\RawBlock}[1]{\par\vspace{3pt}\noindent#1\par\vspace{3pt}}
\newcommand{\CenterLine}[1]{\makebox[\linewidth][c]{#1}}
\newlength{\FullSourceWidth}\newlength{\GutterOffset}\newif\ifNumbersLeft
\newcommand{\DSource}[2]{\BeginSource{AB01-PDF#1}{#2}%
 \setlength{\FullSourceWidth}{\textwidth}\setlength{\GutterOffset}{0pt}%
 \ifodd#2\relax\NumbersLefttrue\else\NumbersLeftfalse\fi}
\newsavebox{\DLBox}
\newcommand{\DLine}[4]{%
 \par\noindent\hypertarget{#1}{}%
 \ifx\relax#2\relax\else
   \ifNumbersLeft
    \rlap{\kern-\GutterOffset\llap{\fontsize{7}{8}\selectfont #2\hspace{3mm}}}%
   \else
    \rlap{\kern\dimexpr\FullSourceWidth-\GutterOffset\relax\hspace{3mm}{\fontsize{7}{8}\selectfont #2}}%
   \fi
 \fi
 \ifx\relax#3\relax\else
   \ifNumbersLeft
    \rlap{\kern\dimexpr\FullSourceWidth-\GutterOffset\relax\hspace{3mm}{\fontsize{7}{8}\selectfont #3}}%
   \else
    \rlap{\kern-\GutterOffset\llap{\fontsize{7}{8}\selectfont #3\hspace{3mm}}}%
   \fi
 \fi
 \sbox{\DLBox}{\strut #4}%
 \ifdim\wd\DLBox>\linewidth\typeout{DIPLOMATICWIDTH|#1|\the\wd\DLBox|\the\linewidth}\fi
 \usebox{\DLBox}\strut\par}
\begin{document}
\thispagestyle{empty}
\begin{center}\Large AL-BATTĀNĪ\par\vspace{4mm}\LARGE OPUS ASTRONOMICUM\par
\vspace{5mm}\large Caroli Alphonsi Nallino versio Latina et adnotationes (pars I)\par\vspace{8mm}
\Large S02 --- editio diplomatica\par\vspace{4mm}\large Paginae impressae 54--55\end{center}
\vspace{12mm}\noindent Singulae lineae paginarum impressarum singulis lineis huius editionis respondent, cum
divisionibus vocabulorum, signis et generibus litterarum (rectis, inclinatis, distantibus). Numeri linearum et
paginae textus Arabici in marginibus stant ut impressi sunt. Textus ex imaginibus paginarum transcriptus est;
litterae Graecae et Arabicae, numeri et signa in imaginibus amplificatis lecta sunt.
Haec transcriptio a nullo homine recognita est.\par
\DSource{0143}{54}
\DLine{AB01-PDF0143-body-L001}{}{}{dum complectatur ex gradibus quorum circumferentia triangulum DKE circumdans 360 habet;}
\DLine{AB01-PDF0143-body-L002}{}{}{eius sinus 0\textsuperscript{\textit{p}} 31′ 25″ (1) erit ex partibus quarum semidiametrus ED 60 continet. Remanent}
\DLine{AB01-PDF0143-body-L003}{}{p. 82.}{igitur angulus KED 89° 30′, arcus KD 89° 30′ (2), eiusque sinus 60\textsuperscript{\textit{p}} (3) fere. Sed si linea ED}
\DLine{AB01-PDF0143-body-L004}{}{}{excentricitatis 10\textsuperscript{\textit{p}} 19′ ponatur, linea EK quinque fere minutorum, et linea KD 10\textsuperscript{\textit{p}} 19′ circiter}
\DLine{AB01-PDF0143-body-L005}{5}{}{fiet. Deinde lineam EF ut ED ponamus, ES ut EK, SF ut DK. Lineam DB, quae a centro}
\DLine{AB01-PDF0143-body-L006}{}{}{excentrici ad eiusdem circumferentiam protenditur, iam vidimus 49\textsuperscript{\textit{p}} 41′ esse ea ratione qua}
\DLine{AB01-PDF0143-body-L007}{}{}{semidiametrus MB epicycli 5\textsuperscript{\textit{p}} 15′ est; ob ea quae diximus tota BK 48\textsuperscript{\textit{p}} 36′ erit, de qua si}
\DLine{AB01-PDF0143-body-L008}{}{}{EK, scilicet 5 minuta, demitur, remanet linea EB 48\textsuperscript{\textit{p}} 31′. Quoniam ES etiam 5 minuta com-}
\DLine{AB01-PDF0143-body-L009}{}{}{plectitur, erit SB 48\textsuperscript{\textit{p}} 26′ (4); et ex ratione linearum SF, SB ratio lineae BF deprehenditur,}
\DLine{AB01-PDF0143-body-L010}{10}{}{quae igitur 49\textsuperscript{\textit{p}} 31′ fere complectetur si SF circiter 10\textsuperscript{\textit{p}} 19′ (5) fuerit. Excrescente autem}
\DLine{AB01-PDF0143-body-L011}{}{}{linea BF usque ad 60\textsuperscript{\textit{p}}, fiet SF circiter 12\textsuperscript{\textit{p}} 32′ (6), et arcus super eam positus, si angulus}
\DLine{AB01-PDF0143-body-L012}{}{}{rectus 90° continet, 12° 1′ fere; et haec est quantitas arcus HG. — Motus verus Lunae in}
\DLine{AB01-PDF0143-body-L013}{}{}{epicyclo, ut a centro eclipticae conspicitur, a puncto G erit igitur 345° 13′; qua re, si dupla}
\DLine{AB01-PDF0143-body-L014}{}{}{elongatio minor est 180° gradibus, arcus HG additur anomaliae in tabulis repertae; si contra}
\DLine{AB01-PDF0143-body-L015}{15}{}{180° superat, arcus HG de anomalia illa demitur. Initio enim centrum epicycli inter pun-}
\DLine{AB01-PDF0143-body-L016}{}{}{ctum A excentrici et P, qui M sequitur, erit; postea in alterum semicirculum PA, qui a}
\DLine{AB01-PDF0143-body-L017}{}{}{parte puncti D reperitur (7), transibit. Haec 12° 1′ in tertia (8) columna, sub numero 90° 30′}
\DLine{AB01-PDF0143-body-L018}{}{}{inveniuntur.}
\hypertarget{AB01-PDF0143-P01}{}
\DLine{AB01-PDF0143-body-L019}{}{}{\Indent Dicit [auctor]: Motum longitudinis Lunae ut in libro Ptolemaei ponitur invenimus, post-}
\DLine{AB01-PDF0143-body-L020}{20}{}{quam ei addidimus id quod motui Solis adiecimus; et hoc modo eum in tabulis descripsimus.}
\DLine{AB01-PDF0143-body-L021}{}{}{Motus anomalisticus omnino est ut apud Ptolemaeum; sed motum latitudinis invenimus 27}
\DLine{AB01-PDF0143-body-L022}{}{}{minutis redundantem in libro Ptolemaei, qua re haec minuta per tempus inter eum et nos}
\DLine{AB01-PDF0143-body-L023}{}{}{transactum divisimus et de motu latitudinis dempsimus; quod remansit in tabulis posuimus.}
\DLine{AB01-PDF0143-body-L024}{}{}{Columnas ad elongationem duplicatam Lunae a Sole designandam nobis componere necesse}
\DLine{AB01-PDF0143-body-L025}{25}{}{non fuit; quocumque tempore enim elongationem Solis a Luna iuxta eorum medium motum}
\DLine{AB01-PDF0143-body-L026}{}{p. 83.}{scimus, qua duplicata utemus pro eo quod ex tabulis provenit, sine augmento vel deminu-}
\DLine{AB01-PDF0143-body-L027}{}{}{tione. — Latitudinem Lunae usque ad 5 gradus fere pervenire reperimus eaque in septima (9)}
\DLine{AB01-PDF0143-body-L028}{}{}{columna tabularum aequationis describitur; sed cum eius quantitas (10) hoc nostro tempore [ab}
\DLine{AB01-PDF0143-body-L029}{}{}{aerae Dhū ’l–qarnayn initio] $\qfrac{3}{4}$ unius gradus [veritatem] excederet, hoc de motu eius de-}
\DLine{AB01-PDF0143-body-L030}{30}{}{traximus, et remansit quod, Deo altissimo volente, eo loco invenies (\textbf{\textit{d}}).}
\hypertarget{AB01-PDF0143-P02}{}
\DLine{AB01-PDF0143-body-L031}{}{}{\Indent Eclipsis lunaris causa haec est: cum Sol maior Terra sit, Terra autem maior quam Luna,}
\DLine{AB01-PDF0143-body-L032}{}{}{radii Solis Terram involvunt et ab altera parte in aerem progrediuntur ut arbor pini, —}
\DLine{AB01-PDF0143-body-L033}{}{}{qua re umbra Terrae \textit{umbra pinea} appellatur, — et ultra sphaeram Mercurii desinit (\textbf{\textit{e}}).}
\DLine{AB01-PDF0143-body-L034}{}{}{Cum Luna per alterutrum nodum suae sphaerae transit tempore quo Soli opponitur, [centrum}
\DLine{AB01-PDF0143-body-L035}{35}{}{eius orbis] fit centrum ipsum eclipticae, qua re [Luna] in zodiacum incidit et Solem de dia-}
\DLine{AB01-PDF0143-body-L036}{}{}{metro sphaerae caelestis prospicit (11); tunc Luna caret latitudine qua a Sole declinet, et dire-}
\par\vspace{6pt}\nointerlineskip\hbox to\textwidth{\hss\vrule width 115.43pt height.25pt depth0pt\hss}\prevdepth=-1000pt \vspace{5pt}
\noindent\begin{minipage}[t]{.48\textwidth}\vspace{0pt}\fontsize{9}{11.5}\selectfont\setlength{\GutterOffset}{0pt}\setlength{\FullSourceWidth}{\linewidth}
\hypertarget{AB01-PDF0143-N01}{}
\DLine{AB01-PDF0143-notes_left-L001}{}{}{\Indent (1) Plato 30.55.}
\hypertarget{AB01-PDF0143-N02}{}
\DLine{AB01-PDF0143-notes_left-L002}{}{}{\Indent (2) Plato semper 88° 30′.}
\hypertarget{AB01-PDF0143-N03}{}
\DLine{AB01-PDF0143-notes_left-L003}{}{}{\Indent (3) In tabulis sinuum 59\textsuperscript{\textit{p}} 59′ 52″.}
\hypertarget{AB01-PDF0143-N04}{}
\DLine{AB01-PDF0143-notes_left-L004}{}{}{\Indent (4) Plat. 48\textsuperscript{\textit{p}} 56′.}
\hypertarget{AB01-PDF0143-N05}{}
\DLine{AB01-PDF0143-notes_left-L005}{}{}{\Indent (5) Plat. 10\textsuperscript{\textit{p}} 18′.}
\hypertarget{AB01-PDF0143-N06}{}
\DLine{AB01-PDF0143-notes_left-L006}{}{}{\Indent (6) Plat. 12\textsuperscript{\textit{p}} 30′.}
\hypertarget{AB01-PDF0143-N07}{}
\DLine{AB01-PDF0143-notes_left-L007}{}{}{\Indent (7) Hae litterae ad figuram pag. 51, non pag. 53,}
\DLine{AB01-PDF0143-notes_left-L008}{}{}{referuntur.}
\hypertarget{AB01-PDF0143-N08}{}
\DLine{AB01-PDF0143-notes_left-L009}{}{}{\Indent (8) Vel potius quarta.}
\hypertarget{AB01-PDF0143-N09}{}
\DLine{AB01-PDF0143-notes_left-L010}{}{}{\Indent (9) Ita quoque in tabulis.}
\end{minipage}
\hfill\vrule width.3pt\hfill\begin{minipage}[t]{.48\textwidth}\vspace{0pt}\fontsize{9}{11.5}\selectfont\setlength{\GutterOffset}{0pt}\setlength{\FullSourceWidth}{\linewidth}
\hypertarget{AB01-PDF0143-N10}{}
\DLine{AB01-PDF0143-notes_right-L001}{}{}{\Indent (10) Cod. « anomalia eius » (\ArabicRun{حاصته}) pro « eius}
\DLine{AB01-PDF0143-notes_right-L002}{}{}{quantitate » (\ArabicRun{حصته}).}
\hypertarget{AB01-PDF0143-N11}{}
\DLine{AB01-PDF0143-notes_right-L003}{}{}{\Indent (11) Cod.: « tempore eius oppositionis contra}
\DLine{AB01-PDF0143-notes_right-L004}{}{}{« Solem fit centrum eclipticae, qua re in zodiacum}
\DLine{AB01-PDF0143-notes_right-L005}{}{}{« incidit »; Plato: « cumque per alterum nodorum}
\DLine{AB01-PDF0143-notes_right-L006}{}{}{« sui caeli praeventionis hora Luna transgreditur,}
\DLine{AB01-PDF0143-notes_right-L007}{}{}{« eius centrum est centrum circuli signorum, et ipsa}
\DLine{AB01-PDF0143-notes_right-L008}{}{}{« tunc in ipso eodem signorum cingulo consistit ».}
\DLine{AB01-PDF0143-notes_right-L009}{}{}{Textus, pag. 83, l. 7 ita restituatur: \ArabicRun{صار مركزه مركز}}
\DLine{AB01-PDF0143-notes_right-L010}{}{}{\ArabicRun{فلك البروج}.}
\end{minipage}
\DSource{0144}{55}
\DLine{AB01-PDF0144-body-L001}{}{}{ctionem Solis quaerit. Terra autem, inter ambo luminaria interiecta, Lunam tegit, quae igitur}
\DLine{AB01-PDF0144-body-L002}{}{}{in umbram pineam Terrae incidit, lumen suum amittit, et plus minusve occultatur prout eius}
\DLine{AB01-PDF0144-body-L003}{}{}{iter [i. e. longitudo] plus vel minus [punctum] zodiaci appropinquat, in quo centrum umbrae}
\DLine{AB01-PDF0144-body-L004}{}{}{reperitur. Si Luna latitudine omnino caret et in ipso nodo exstat, tempore mediae eclipsis}
\DLine{AB01-PDF0144-body-L005}{5}{}{per lineam centri umbrae transit; eclipsis igitur omnium perfectissima ac longissima est, et}
\DLine{AB01-PDF0144-body-L006}{}{}{Luna totum suum lumen amittit. Nam tota lux profecto tum habetur cum Sol super totum}
\DLine{AB01-PDF0144-body-L007}{}{}{hemisphaerium lunare eum prospiciens incidit, quod fieri non potest nisi (1) Luna in diametro}
\DLine{AB01-PDF0144-body-L008}{}{}{Solis sit, ita ut a Sole dimidia circumferentia distet; tunc vero Luna, [cum res sint ut ita}
\DLine{AB01-PDF0144-body-L009}{}{}{diximus], est sub umbra in medio eclipsis. Sed si a via Solis in septentriones aut meridiem}
\DLine{AB01-PDF0144-body-L010}{10}{}{versus declinat, inter ambo luminaria distantia semicirculo minor habetur, neque oppositio}
\DLine{AB01-PDF0144-body-L011}{}{}{iuxta diametrum rectam fit; qua re cum Luna opponitur Soli ab eius via in latitudine decli-}
\DLine{AB01-PDF0144-body-L012}{}{}{nans, eclipsis fit iuxta quantitatem huius latitudinis, usque ad latitudinem maximam cuius cir-}
\DLine{AB01-PDF0144-body-L013}{}{}{culus circulum umbrae tangere possit.}
\hypertarget{AB01-PDF0144-P01}{}
\DLine{AB01-PDF0144-body-L014}{}{}{\Indent Ex his patet reliquos planetas in oppositione Solis occultari non posse, nam umbra ad}
\DLine{AB01-PDF0144-body-L015}{15}{}{eorum centra non pervenit; Mercurius autem a Sole nunquam ita recedit ut ei opponatur (\textbf{\textit{f}}),}
\DLine{AB01-PDF0144-body-L016}{}{}{qua re in umbram de oppositione non incidit. Sed Luna reliquos planetas occultat, quando in}
\DLine{AB01-PDF0144-body-L017}{}{}{longitudine et latitudine eius viam apparentem occupant; tunc enim Luna in lineam ab oculis}
\DLine{AB01-PDF0144-body-L018}{}{}{ad planetam progredientem incidit. Planetae interdum alios [planetas] occultant, cum inferior}
\DLine{AB01-PDF0144-body-L019}{}{p. 84.}{eorum sub superiore transeat eadem longitudine et latitudine.}
\hypertarget{AB01-PDF0144-P02}{}
\DLine{AB01-PDF0144-body-L020}{20}{}{\Indent Si Terra Solem magnitudine aequaret, latitudo umbrae sibi ipsi aequalis ubique maneret,}
\DLine{AB01-PDF0144-body-L021}{}{}{neque in aere terminaretur, sed ad infinitum procederet; eclipsis lunaris (2) semper eadem}
\DLine{AB01-PDF0144-body-L022}{}{}{quantitate sive in inferiore sive in superiore parte epicycli Lunae eveniret; planetae autem}
\DLine{AB01-PDF0144-body-L023}{}{}{omnes (3) in oppositione Solis occultarentur. Si Sol minor quam Terra esset, summa umbra}
\DLine{AB01-PDF0144-body-L024}{}{}{latior inferiore exstaret, et in Caelo ad infinitum procederet; quo magis latitudo umbrae}
\DLine{AB01-PDF0144-body-L025}{25}{}{ascenderet, eo magis Luna et planetae occultarentur, atque diversis diebus deficerent iuxta}
\DLine{AB01-PDF0144-body-L026}{}{}{suum iter ab itinere Solis.}
\hypertarget{AB01-PDF0144-P03}{}
\DLine{AB01-PDF0144-body-L027}{}{}{\Indent Causa eclipsium solarium est Luna; haec enim, si contingit ut eius centrum tempore}
\DLine{AB01-PDF0144-body-L028}{}{}{coniunctionis appareat in ecliptica, inter oculos adspicientium et Solem incidit, quia est in linea}
\DLine{AB01-PDF0144-body-L029}{}{}{ab oculis ad Solem ducta. Eo tempore Luna est Terrae propior quam Sol, et similis rei parvae}
\DLine{AB01-PDF0144-body-L030}{30}{}{quae semper videatur atque alteram maiorem occultet, si magis quam haec propinqua sit}
\DLine{AB01-PDF0144-body-L031}{}{}{oculis. Quantitates eclipsis variant iuxta quantitatem latitudinis apparentis Lunae, donec [lati-}
\DLine{AB01-PDF0144-body-L032}{}{}{tudo] talis fiat ut nihil Solis occultare queat. Eclipsis Solis diversam quantitatem habet in}
\DLine{AB01-PDF0144-body-L033}{}{}{locis [Terrae] qui diversam ab alteris latitudinem habent; eclipsis contra Lunae ubique aequa-}
\DLine{AB01-PDF0144-body-L034}{}{}{lis apparet.}
\hypertarget{AB01-PDF0144-P04}{}
\DLine{AB01-PDF0144-body-L035}{35}{}{\Indent Ut cognosceret Solis et Lunae distantias diametrosque, et eorum corporum respectu Terrae}
\DLine{AB01-PDF0144-body-L036}{}{}{magnitudines, Ptolemaeus (4) duabus eclipsibus Lunae usus est; supposuitque, ad calculos}
\DLine{AB01-PDF0144-body-L037}{}{}{ducendos, Lunam totum Solem occultare posse, cum Luna suum apogeum tempore coniun-}
\DLine{AB01-PDF0144-body-L038}{}{}{ctionis occupet et in ecliptica consistat; nec diametro Solis in apogeo et perigeo differentiam}
\DLine{AB01-PDF0144-body-L039}{}{}{sensibilem respectu Lunae tribuit, immo diametrum semper aequalem esse posuit. Denique in}
\par\vspace{6pt}\nointerlineskip\hbox to\textwidth{\hss\vrule width 115.43pt height.25pt depth0pt\hss}\prevdepth=-1000pt \vspace{5pt}
\noindent\begin{minipage}[t]{.48\textwidth}\vspace{0pt}\fontsize{9}{11.5}\selectfont\setlength{\GutterOffset}{0pt}\setlength{\FullSourceWidth}{\linewidth}
\hypertarget{AB01-PDF0144-N01}{}
\DLine{AB01-PDF0144-notes_left-L001}{}{}{\Indent (1) Plato: « nam totum lumen non amittet, nisi}
\DLine{AB01-PDF0144-notes_left-L002}{}{}{« cum super illam sui circuli partem qui in ipsius}
\DLine{AB01-PDF0144-notes_left-L003}{}{}{« directo fuerit, ceciderit, quod esse non potest,}
\DLine{AB01-PDF0144-notes_left-L004}{}{}{« nisi ..... ».}
\hypertarget{AB01-PDF0144-N02}{}
\DLine{AB01-PDF0144-notes_left-L005}{}{}{\Indent (2) Addendum est « perfecta »; sin minus haud}
\DLine{AB01-PDF0144-notes_left-L006}{}{}{verum est.}
\end{minipage}
\hfill\vrule width.3pt\hfill\begin{minipage}[t]{.48\textwidth}\vspace{0pt}\fontsize{9}{11.5}\selectfont\setlength{\GutterOffset}{0pt}\setlength{\FullSourceWidth}{\linewidth}
\hypertarget{AB01-PDF0144-N03}{}
\DLine{AB01-PDF0144-notes_right-L001}{}{}{\Indent (3) Mercurio et Venere tamen exceptis, qui}
\DLine{AB01-PDF0144-notes_right-L002}{}{}{nunquam, ut supra auctor de Mercurio dixit, Soli}
\DLine{AB01-PDF0144-notes_right-L003}{}{}{opponuntur. Addendum ergo videtur « superiores ».}
\hypertarget{AB01-PDF0144-N04}{}
\DLine{AB01-PDF0144-notes_right-L004}{}{}{\Indent (4) \textit{Almag.} V, 14 (ed. Halma, t. I, p. 339-343).}
\DLine{AB01-PDF0144-notes_right-L005}{}{}{De iis cfr. \Name{Ginzel}, \textit{Spezieller Kanon der Sonnen-}}
\DLine{AB01-PDF0144-notes_right-L006}{}{}{\textit{und Mondfinsternisse,} Berlin 1899, V Abschnitt.}
\end{minipage}
\end{document}
